\section{Discussion}

The experiments distinguish two mechanisms that are often combined inside a
single learned dynamics model. Complementary context correction reduces the
forecast-origin error in (\ref{eq:rollout_bound}); the learned residual is
intended to reduce accumulated transition error. Here the first mechanism is
consistently useful under observation corruption. It explains the
\CorruptImprovement\% open-loop gain, the outage response, and the reduction in
closed-loop action effort relative to PGNN-style. Its small closed-loop state
loss advantage is repeated across all system--condition means, although the
aggregate magnitude is only \ClosedLoopImprovement\%.

The residual evidence points in the other direction. Filter plus nominal
dynamics is better than the full model in both the post-hoc open-loop analysis
and the frozen closed-loop study. This result is compatible with a nominal
model that is accurate enough that the learned correction adds estimation
variance or compounds small one-step biases. Residual learning may still be
useful when mismatch is larger, but the current benchmark does not demonstrate
that benefit. A model-selection rule for automation should consequently retain
the zero-residual variant rather than assuming that an additional trainable
component improves control.

Equation~(\ref{eq:alpha}) also explains the clean-regime result. A fixed gain
of 0.5 assumes a particular relationship between nominal-prediction and
measurement errors. When measurements are clean, equal weighting can discard
useful measurement information; when measurements are corrupted, it can
suppress noise and bridge missing channels. A channel- and state-dependent
gain estimated without access to test outcomes is a direct extension. Its
evaluation should retain the same clean noninferiority gate because an
adaptive estimator can otherwise trade nominal accuracy for corruption
robustness.

The measured-flight study narrows the transfer claim. Position error improves,
whereas full-state error and interval coverage do not. The reduced quadrotor
prior omits platform effects, and the diagonal Gaussian head does not represent
cross-state or temporal dependence. Risk-sensitive control therefore requires
recalibration or a richer predictive distribution before using these
uncertainties as constraint margins.

The closed-loop experiment uses simulated plants, one deterministic-mean CEM
configuration, 20 control steps, and a quadratic regulation objective. The
measured flights provide real sensor and actuation records but not an online
intervention. Physical closed-loop experiments, task constraints, longer
missing-data bursts, informative sensor failures, and latency on the target
controller remain necessary to assess deployment performance. The bounded
residual and finite-horizon results do not by themselves establish closed-loop
stability.

\section{Conclusion}

CFR-WM separates a mask-gated complementary state update from a
nominal-plus-residual free rollout. Across a controlled open-loop benchmark,
the interface improves 25-step prediction under corrupted observations but
does not meet the clean-regime noninferiority criterion. In a frozen
receding-horizon study it reduces state loss relative to PGNN-style by
\ClosedLoopImprovement\%, while filter plus nominal dynamics remains slightly
better than the complete model. Measured flights show a position benefit
alongside worse full-state error and undercovered intervals. For the tested
automation setting, complementary context filtering is the supported design
element; residual and uncertainty components require stronger evidence before
they are used for control decisions.
